% models/mathematical/quantenzeit_operatoren.tex
% Quantenzeit: Operatoren, Emergenz und Korrelationen

\section{Quantenzeit und Operatorstruktur}

\subsection{Zeit als nicht-existenter Operator}

In der Standard-Quantenmechanik existiert kein Zeitoperator:
\[
\hat{T} \notin \mathcal{H}.
\]

Zeit tritt als Parameter $t$ auf, nicht als Observablen-Operator.

Dies motiviert die Interpretation von Zeit als emergente Größe.

\subsection{Page--Wootters-Formalismus}

Zeit entsteht aus Korrelationen zwischen System und Uhr:
\[
|\Psi\rangle = \sum_t |t\rangle_{\text{Uhr}} \otimes |\psi(t)\rangle_{\text{System}}.
\]

Die Uhr definiert eine bedingte Dynamik:
\[
\hat{H} |\Psi\rangle = 0.
\]

\subsection{Entanglement als Zeitgenerator}

Die fundamentale Zeitrichtung ergibt sich aus:
\[
\frac{d}{dt} \rho = \text{Tr}_{\text{Uhr}} \big[ \hat{H}, \rho_{\text{gesamt}} \big].
\]

Zeit ist damit ein Maß für Entanglement-Fluss.

\subsection{Majorana-Symmetrien}

Reelle Zustände erfüllen:
\[
\psi = \psi^*.
\]

Die Zeitumkehrsymmetrie $T$ wirkt als:
\[
T \psi(t) = \psi(-t).
\]

Majorana-Zustände definieren eine fundamentale Zeitrichtung durch
Symmetriebrechung:
\[
T^2 = +1.
\]

\subsection{Quantenzeit-Operatoren (effektiv)}

Obwohl kein fundamentaler Zeitoperator existiert, lassen sich effektive
Operatoren definieren:

\[
\hat{\Theta} = i \frac{\partial}{\partial E},
\quad
\hat{\Omega} = -i \frac{\partial}{\partial t}.
\]

Diese Operatoren beschreiben Energie-Zeit-Dualität.

\subsection{Emergente Zeitrate}

Die effektive Zeitrate ergibt sich aus:
\[
\frac{d\tau}{dt} = \frac{\partial S}{\partial E}
\]
mit der Wirkung $S$.

\subsection{Kopplung an das Zeitfeld}

Die quantenphysikalische Schicht trägt zum Zeitfeldtensor bei durch:
\[
\mathcal{T}^{(q)}_{00} = \langle \Psi | \hat{H} | \Psi \rangle,
\]
\[
\mathcal{T}^{(q)}_{0i} = \langle \Psi | \hat{J}_i | \Psi \rangle,
\]
\[
\mathcal{T}^{(q)}_{ij} = C_{ij}(\Psi),
\]
wobei $C_{ij}$ die Korrelationsmatrix ist.

\subsection{Interferenz mit höheren Schichten}

Die Kopplung erfolgt über:
\[
\mathcal{K}\big(\mathcal{T}^{(q)}, \mathcal{T}^{(i)}\big)
= \lambda_{qi} \, C_{ij} \, \mathcal{T}^{(i)\,ij}.
\]

\subsection{Ergebnis}

Die quantenphysikalische Schicht definiert:

\begin{itemize}
\item die fundamentale Zeitrichtung,
\item die emergente Dynamik,
\item die Korrelationen, die höhere Schichten strukturieren.
\end{itemize}

Sie bildet die Basis des gesamten Zeitfeldes.
